% Wygenerowano z ../06_detektor.md skryptem _md_do_latex.py (można edytować ręcznie).
\section{Ćwiczenie 6: gotowa sieć neuronowa szuka obiektów, czyli jak zapytać sieć (20 min)}\label{ux107wiczenie-6-gotowa-sieux107-neuronowa-szuka-obiektuxf3w-czyli-jak-zapytaux107-sieux107-20-min}

\textbf{Notatnik:} \texttt{06\_detektor.ipynb}

\textbf{Potrzebne:} obiekt do szukania (np. czerwona gumowa kaczka) na podłodze przed robotem, kilka innych przedmiotów (krzesło, plecak, pudełko), miarka, „obiekty-zagadki'' od prowadzącego

\subsection{Cel}\label{cel}

Sprawdzić, co potrafią gotowe sieci rozpoznające obiekty, nauczyć się dobierać zapytanie słowne i wybrać ustawienia do zadania „znajdź obiekt'' w ćwiczeniu 7.

\subsection{Trochę teorii}\label{trochux119-teorii}

Porównamy dwa rodzaje gotowych, wytrenowanych sieci:

\begin{itemize}
\tightlist
\item
  \textbf{Sieć z zamkniętą listą klas (YOLO, zbiór COCO).} Nauczona na ponad 100 000 opisanych zdjęć rozpoznaje \textbf{80 rodzajów obiektów}: osobę, krzesło, plecak, butelkę, kubek i inne. Czego nie ma na liście, tego nie wykryje. Najwyżej pomyli to z czymś podobnym.
\item
  \textbf{Sieć z otwartym słownikiem (YOLOE).} Tej sieci mówisz słowami (po angielsku), czego ma szukać, np. \texttt{"rubber\ duck"}. Nauczyła się łączyć obrazy z opisami słownymi, więc znajduje obiekty, na których \textbf{nikt jej specjalnie nie uczył}. Wynik zależy jednak od tego, \textbf{jak zapytasz}.
\end{itemize}

Dla każdego wykrytego obiektu sieć podaje \textbf{ramkę} (gdzie on jest) i \textbf{pewność} (od 0 do 1). Wykrycia słabsze od \textbf{progu pewności} są odrzucane.

Sieci działają na \textbf{procesorze} stacji. Jedno zdjęcie analizują przez kilkaset milisekund, czyli ok. 2--10 zdjęć na sekundę.

\subsection{Kroki}\label{kroki}

\subsubsection{Krok 1: sieć z zamkniętą listą (3 min)}\label{krok-1-sieux107-z-zamkniux119tux105-listux105-3-min}

Postaw kaczkę ok. 1 m przed robotem, a obok krzesło i plecak. Uruchom komórkę z modelem \texttt{"coco"}.

\begin{itemize}
\tightlist
\item
  Co wykryła sieć? \linia{14.0}
\item
  Czy wykryła kaczkę? Jeśli tak, to jako co? \linia{9.2}
\end{itemize}

\subsubsection{Krok 2: zapytaj słowami (5 min)}\label{krok-2-zapytaj-sux142owami-5-min}

Uruchom komórkę \texttt{porownaj\_slowa}. Sprawdza ona wiele opisów na tym samym zdjęciu i pokazuje, przy którym sieć jest najpewniejsza. Dopisz do listy \textbf{własne pomysły} (co najmniej 4).

\par\noindent\begin{minipage}{\linewidth}
\begin{tabular}{@{}
  >{\raggedright\arraybackslash}p{(\linewidth - 4\tabcolsep) * \real{0.3846}}
  >{\raggedright\arraybackslash}p{(\linewidth - 4\tabcolsep) * \real{0.3077}}
  >{\raggedright\arraybackslash}p{(\linewidth - 4\tabcolsep) * \real{0.3077}}@{}}
\toprule\noalign{}
\begin{minipage}[b]{\linewidth}\raggedright
Opis (zapytanie)
\end{minipage} & \begin{minipage}[b]{\linewidth}\raggedright
Pewność
\end{minipage} & \begin{minipage}[b]{\linewidth}\raggedright
Uwagi
\end{minipage} \\
\midrule\noalign{}
\texttt{duck} & & \\
\texttt{rubber\ duck} & & \\
\texttt{kaczka} (po polsku) & & \\
\texttt{bird} & & \\
& & \\
& & \\
& & \\
& & \\
\bottomrule\noalign{}
\end{tabular}
\end{minipage}\par\medskip

Najlepsze zapytanie: \linia{9.2}

\subsubsection{Krok 3: słowa konkurują (2 min)}\label{krok-3-sux142owa-konkurujux105-2-min}

Ustaw w \texttt{ustaw\_slowa} jednocześnie Twoje najlepsze zapytanie i inne słowa, np. \texttt{{[}"rubber\ duck",\ "box",\ "toy"{]}}. Czy pewność wykrycia kaczki się zmieniła? Czy pojawiły się nowe wykrycia? \linia{9.2}

\subsubsection{Krok 4: obiekt-zagadka (4 min)}\label{krok-4-obiekt-zagadka-4-min}

Prowadzący daje zespołowi przedmiot. Ustawcie go przed robotem i znajdźcie zapytanie, przy którym sieć wykrywa go z pewnością \textbf{powyżej 0,5}. Wygrywa zespół, który zrobi to najszybciej.

Przedmiot: \linia{6.7} najlepsze zapytanie: \linia{6.7} pewność: \linia{2.5}

\subsubsection{Krok 5: odległość, rozdzielczość, szybkość (4 min)}\label{krok-5-odlegux142oux15bux107-rozdzielczoux15bux107-szybkoux15bux107-4-min}

Ustawiaj kaczkę w różnych odległościach. Dla każdej odległości uruchom komórkę, która sprawdza rozdzielczości 320 i 640.

\par\noindent\begin{minipage}{\linewidth}
\begin{tabular}{@{}
  >{\raggedright\arraybackslash}p{(\linewidth - 4\tabcolsep) * \real{0.2549}}
  >{\raggedright\arraybackslash}p{(\linewidth - 4\tabcolsep) * \real{0.3725}}
  >{\raggedright\arraybackslash}p{(\linewidth - 4\tabcolsep) * \real{0.3725}}@{}}
\toprule\noalign{}
\begin{minipage}[b]{\linewidth}\raggedright
Odległość {[}m{]}
\end{minipage} & \begin{minipage}[b]{\linewidth}\raggedright
Pewność przy 320 px
\end{minipage} & \begin{minipage}[b]{\linewidth}\raggedright
Pewność przy 640 px
\end{minipage} \\
\midrule\noalign{}
0,5 & & \\
1,0 & & \\
2,0 & & \\
3,0 & & \\
\bottomrule\noalign{}
\end{tabular}
\end{minipage}\par\medskip

Szybkość: 320 px: \linia{2.5} zdjęć/s, 640 px: \linia{2.5} zdjęć/s

\subsubsection{Krok 6: decyzja zespołu (2 min)}\label{krok-6-decyzja-zespoux142u-2-min}

Wybierzcie ustawienia do ćwiczenia 7, w którym robot ma sam znaleźć obiekt i do niego podjechać.

\begin{itemize}
\tightlist
\item
  Zapytanie: \linia{6.7} próg: \linia{2.5} rozdzielczość: \linia{2.5}
\item
  Uzasadnienie: \linia{14.0}
\end{itemize}

\subsection{Pytania do przemyślenia}\label{pytania-do-przemyux15blenia}

\begin{enumerate}
\def\labelenumi{\arabic{enumi}.}
\tightlist
\item
  Dlaczego sieć z zamkniętą listą nie może wykryć kaczki, choć „widzi'' ją na zdjęciu?
\item
  Dlaczego zapytanie po polsku działa gorzej niż po angielsku? \emph{(Podpowiedź: na jakich opisach uczono sieć?)}
\item
  Robot ma szukać obiektu w całym pomieszczeniu. Co jest ważniejsze: duża rozdzielczość (widać daleko) czy szybkość? Czy da się połączyć jedno i drugie, np. szukać dokładnie, a podjeżdżać szybko?
\item
  Co jest gorsze dla robota szukającego obiektu: \textbf{fałszywe wykrycie} (pojedzie do czegoś innego) czy \textbf{przeoczenie}? Jak to wpływa na wybór progu?
\end{enumerate}
